Recall \wikilink{s13--b008}{the derived series} and
\wikilink{s13--b010}{solvability}.

\begin{proposition}\label{derived-homomorphisms}
For a surjective homomorphism $\phi:G\to Q$, one has
$\phi(D^r(G))=D^r(Q)$ for every term of the derived series, where
$D^0(G)=G$ and $D^{r+1}(G)=[D^r(G),D^r(G)]$.
\end{proposition}
\begin{proof}\label{derived-homomorphisms-proof}
Homomorphisms send commutators to the commutators of their images.
Surjectivity supplies lifts of all pairs of elements of $Q$, so the
images of the generated commutator subgroups are equal. Repeat this
argument on each derived subgroup.
\end{proof}
\begin{theorem}\label{solvable-extensions}
If $N\trianglelefteq G$ and both $N$ and $G/N$ are solvable, then
$G$ is solvable. No finiteness assumption is needed.
\end{theorem}
\begin{proof}\label{solvable-extensions-proof}
If the quotient's derived series vanishes at step $r$, the proposition
applied to $G\to G/N$ gives $D^r(G)\subseteq N$.
Commutator formation preserves inclusion, so
$D^{r+s}(G)\subseteq D^s(N)$ for every $s$. Choose $s$ for which
$D^s(N)=1$ to conclude.
\end{proof}
\begin{remark}\label{derived-notation}
The notation $D^r(G)$ distinguishes the derived series from the lower
central series. The proof also bounds the derived length of $G$ by
the sum of the derived lengths of $N$ and $G/N$.
\end{remark}
